(1) 两个自旋耦合满足 \(\mathbf S_1\cdot\mathbf S_2=[(\mathbf S_1+\mathbf S_2)^2-S_1^2-S_2^2]/2\)。单重态总自旋0，能量 \(E_s=-3J/4\)；三重态总自旋1，能量 \(E_t=J/4\)，简并度3。故 \(Z=e^{3\beta J/4}+3e^{-\beta J/4}\)。

(2) 热态是 \(\rho=p_s|s\rangle\langle s|+(1-p_s)\Pi_t/3\)，其中 \(p_s=e^{\beta J}/(e^{\beta J}+3)\)。\(T\to0\) 时 \(p_s\to1\)，为纯自旋单重态；高温 \(\beta J\to0\) 时 \(p_s\to1/4\)，为四维单位混合态。

(3) 在计算基 \(|\uparrow\uparrow\rangle,|\uparrow\downarrow\rangle,|\downarrow\uparrow\rangle,|\downarrow\downarrow\rangle\) 中展开 \(\rho\) 并对第二自旋转置，部分转置谱为三个相等的正本征值 \((1+2p_s)/6\) 以及一个 \(\boxed{\lambda_{\min}=1/2-p_s}\)。因此纠缠当且仅当 \(p_s>1/2\)，即 \(e^{\beta J}>3\)，得到 \(\boxed{T<T_c=J/(k_B\ln3)}\)；该两量子比特情形部分转置判据充要。

(4) 磁场仅改变三重态 \(m=1,0,-1\) 的能量为 \(J/4-m\mu B\)，单重态不变。故 \(Z(B)=e^{3\beta J/4}+e^{-\beta J/4}[1+2\cosh(\beta\mu B)]\)。取 \(M=(1/\beta)\partial_B\ln Z\) 并在 \(B=0\) 求导，得到 \(\boxed{\chi=2\beta\mu^2/(e^{\beta J}+3)}\)，低温因单重态与磁性三重态有能隙 \(J\) 而指数冻结。
